% TOP_PROBLEM: {"id": 30002146, "title": "Nonaxisymmetric Stability of Extremal Kerr Waves", "category_id": 16, "category": {"id": 16, "name": "physics", "display_name": "Mathematical Physics", "description": "Problems at the intersection of mathematics and physics.", "slug": "physics", "order_index": 16, "created_at": "2026-07-31T15:26:25.671Z"}, "status": "open", "statusQualification": "The decay rate was recovered from the primary report (PDF page 30); the data norm and time-foliation conventions remain those of that estimate. Its target is decay of the scalar itself, and known transverse-derivative instabilities are a separate phenomenon."}
% FIRST_POSTED: 2026-10-11
% ENTRY_KIND: statement_only
% SUBMISSION: {"kind": "statement_only", "category": "top_problems", "problem_id": "30002146", "proposal_type": "established", "source": "Mathematical Aspects of General Relativity (2012), Oberwolfach Reports, DOI 10.4171/owr/2012/37, pp. 2269--2333. Primary report accessed 2026-10-11; relevant PDF pages 30, 31. https://doi.org/10.4171/owr/2012/37; Aretakis, axisymmetric waves on extreme Kerr (catalogue-listed context; not independently consulted in this preparation). https://arxiv.org/abs/1110.2006; Dafermos--Rodnianski--Shlapentokh-Rothman, subextremal Kerr stability (catalogue-listed context; not independently consulted in this preparation). https://arxiv.org/abs/1402.7034", "status_evidence": "Canonical source status open retained. Supplied literature triage dated 2026-08-21: Axisymmetric extremal-Kerr decay and non-extremal non-axisymmetric stability are known, but the full non-axisymmetric extremal-Kerr stability problem identified in the record remains unresolved. Primary-report access checked 2026-10-11; this historical report and inherited triage do not certify current resolution of every component.", "status_checked_on": "2026-10-11", "duplicate_check": "Compared canonical ID, title and normalized target with the complete local UnsolvedMath registry; normalized exact-text matches: none; title matches: none. Related or overlapping canonical records: none identified in this batch. Primary source context and the existing OWR statement audit used for scope and duplicate review; an exhaustive mathematical-equivalence/literature search was not performed.", "attribution": "Original proposer not named in the supplied canonical record; attribution remains to the identified Oberwolfach contribution/report, not its editors as invented proposers. Explanatory statement preparation and source/semantic audit assisted by OpenAI GPT-6 Codex; no human mathematical authorship or independent proof verification is claimed."}
% SOURCE_STATEMENT: "For sufficiently regular solutions $\\psi$ of the scalar wave equation on extremal Kerr, does the pointwise decay estimate up to and including the event horizon that is known for axisymmetric solutions remain valid without assuming axisymmetry?"
% Preparation and final self-audit: CONTRIBUTING.md prompts C and D; source/audit artifacts retained in the local physics_statements_2026_10_11 batch.
% Imported from: physics_statements.tex; UnsolvedMath ID 30002146
% !TeX program = lualatex
\documentclass[11pt]{article}
\usepackage[margin=25mm]{geometry}
\usepackage{fontspec}
\setmainfont{Latin Modern Roman}
\usepackage{amsmath,amssymb,mathtools}
\usepackage{unicode-math}
\setmathfont{Latin Modern Math}
\usepackage{xurl,hyperref}
\hypersetup{colorlinks=true,urlcolor=blue}
\setcounter{secnumdepth}{0}
\setlength{\parindent}{0pt}
\setlength{\parskip}{5pt}
\setlength{\emergencystretch}{3em}
\newcommand{\sref}[2]{\hyperlink{source:#1:#2}{[S#2]}}
\newcommand{\cataloguescope}[1]{\paragraph{Recorded target.}#1}
\newcommand{\R}{\mathbb R}
\newcommand{\C}{\mathbb C}
\newcommand{\N}{\mathbb N}
\newcommand{\Z}{\mathbb Z}
\newcommand{\Q}{\mathbb Q}
\newcommand{\F}{\mathbb F}
\newcommand{\eps}{\varepsilon}
\begin{document}
% UnsolvedMath ID 30002146; source catalogue code OWR-12009-002
\setcounter{section}{30002145}
\section{Nonaxisymmetric Stability of Extremal Kerr Waves}\label{30002146}
{\small UnsolvedMath ID 30002146 · Mathematical Physics}
\subsection{Definitions and mathematical statement}

\paragraph{Shared notation.}$\mathbb N=\{1,2,\ldots\}$; the standard number fields have their usual meanings. All additional conventions are stated locally.


\paragraph{Definitions and conventions.}
For the extremal Kerr metric $g_{M,a}$ with $M>0$ and $|a|=M$, the scalar wave equation is $\Box_g\psi=0$, where $\Box_g$ is the metric wave operator. Axisymmetry means invariance under the rotational Killing field $\Phi$, namely $\Phi\psi=0$. The event horizon is the boundary of the region that can communicate with future null infinity. Pointwise decay here concerns $\psi$ itself on the exterior together with the horizon, rather than boundedness or decay of every transverse derivative. The source estimate is $|\psi(\tau,\cdot)|\leq CD\,\tau^{-3/5}$ up to and including the horizon, with $D$ controlled by its sufficiently regular initial data. The question asks to remove axisymmetry from this estimate.

\paragraph{Target retained from the canonical source catalogue.}
\paragraph{Catalogue wording retained for traceability.}
For sufficiently regular solutions $\psi$ of the scalar wave equation on extremal Kerr, does the pointwise decay estimate up to and including the event horizon that is known for axisymmetric solutions remain valid without assuming axisymmetry?

\paragraph{Source scope and qualification.}
The decay rate was recovered from the primary report (PDF page 30); the data norm and time-foliation conventions remain those of that estimate. Its target is decay of the scalar itself, and known transverse-derivative instabilities are a separate phenomenon.
\subsection{Short English statement}
Does the known horizon-inclusive decay of axisymmetric waves on an extremal Kerr black hole extend to waves without rotational symmetry?
\subsection{Sources}
\begin{itemize}
\item[\hypertarget{source:30002146:1}{S1}] ULAM.AI and the UnsolvedMath Contributors, UnsolvedMath, record 30002146 (OWR-12009-002). Catalogue curation under CC BY 4.0; definitions and English exposition added locally.
\url{https://huggingface.co/datasets/ulamai/UnsolvedMath}.
\item[\hypertarget{source:30002146:2}{S2}] Mathematical Aspects of General Relativity (2012), Oberwolfach Reports, DOI 10.4171/owr/2012/37, pp. 2269--2333. Primary report accessed 2026-10-11; relevant PDF pages 30, 31.
\url{https://doi.org/10.4171/owr/2012/37}.
\item[\hypertarget{source:30002146:3}{S3}] Aretakis, axisymmetric waves on extreme Kerr (catalogue-listed context; not independently consulted in this preparation).
\url{https://arxiv.org/abs/1110.2006}.
\item[\hypertarget{source:30002146:4}{S4}] Dafermos--Rodnianski--Shlapentokh-Rothman, subextremal Kerr stability (catalogue-listed context; not independently consulted in this preparation).
\url{https://arxiv.org/abs/1402.7034}.
\end{itemize}
\label{end:30002146}
\end{document}
